%==================================================
% 1. 1.3 冒頭のスローガン差し替え
%==================================================




%==================================================
% 2--3. Example 2 の差し替え
%==================================================




%==================================================
% 4,5 は削除済み
%==================================================


%==================================================
% 6. Example 7 の差し替え
%==================================================



%==================================================
% 7. Theorem 8 の差し替え
%==================================================

\begin{theorem}[区分求積法]
関数 $f:[a,b]\to\mathbb R$ は Riemann 可積分とする．
区間 $[a,b]$ の分割と評価点の列
\[
\mathcal P_k:\ a=x_0^{(k)}<x_1^{(k)}<\cdots<x_{n_k}^{(k)}=b,
\qquad
x_{i-1}^{(k)}\le t_i^{(k)}\le x_i^{(k)}
\]
をとり，
\[
|\mathcal P_k|\to 0\qquad (k\to\infty)
\]
とする．このとき，
\[
S(f,\mathcal P_k,t^{(k)})\to \int_a^b f(x)\,dx
\qquad (k\to\infty)
\]
が成り立つ．
\end{theorem}

\begin{proof}
\[
I:=\int_a^b f(x)\,dx,
\qquad
M:=\sup_{x\in[a,b]}|f(x)|<\infty
\]
とおく．

まず
\[
U(f,\mathcal P_k)-L(f,\mathcal P_k)\to 0
\qquad (k\to\infty)
\]
を示す．
任意の $\varepsilon>0$ をとる．
$f$ は Darboux 可積分であるから，ある分割
\[
\mathcal Q:\ a=y_0<y_1<\cdots<y_m=b
\]
が存在して
\[
U(f,\mathcal Q)-L(f,\mathcal Q)<\frac{\varepsilon}{3}
\]
となる．

各 $k$ について，共通細分
\[
\mathcal R_k:=\mathcal P_k\cup\mathcal Q
\]
を考える．
$\mathcal R_k$ は $\mathcal Q$ の細分だから
\[
U(f,\mathcal R_k)-L(f,\mathcal R_k)
\le
U(f,\mathcal Q)-L(f,\mathcal Q)
<
\frac{\varepsilon}{3}.
\]

一方，$\mathcal R_k$ は $\mathcal P_k$ に $\mathcal Q$ の分点を高々 $m-1$ 個付け加えたものである．
分割に 1 点を付け加えるとき，上和の減少量も下和の増加量も高々 $2M|\mathcal P_k|$ である
（実際，長さ $\ell\le |\mathcal P_k|$ の 1 つの区間を 2 つに分けるだけなので，
旧寄与と新寄与の差は高々 $2M\ell$ で抑えられる）．
したがって
\[
0\le U(f,\mathcal P_k)-U(f,\mathcal R_k)\le 2M(m-1)|\mathcal P_k|,
\]
\[
0\le L(f,\mathcal R_k)-L(f,\mathcal P_k)\le 2M(m-1)|\mathcal P_k|.
\]
$|\mathcal P_k|\to 0$ より，十分大きい $k$ に対して両方とも $\varepsilon/3$ 未満となる．
このとき
\[
\begin{aligned}
U(f,\mathcal P_k)-L(f,\mathcal P_k)
&=
\bigl(U(f,\mathcal P_k)-U(f,\mathcal R_k)\bigr)
+\bigl(U(f,\mathcal R_k)-L(f,\mathcal R_k)\bigr)\\
&\qquad
+\bigl(L(f,\mathcal R_k)-L(f,\mathcal P_k)\bigr)
<\varepsilon.
\end{aligned}
\]
よって
\[
U(f,\mathcal P_k)-L(f,\mathcal P_k)\to 0.
\]
\end{proof}


%==================================================
% 8. Theorem 9 の差し替え
%==================================================

\begin{theorem}[微分積分学の第二基本定理]
$C^1$ 級関数 $F:[a,b]\to\mathbb R$ に対して，$f:=F'$ は区間 $[a,b]$ 上で Riemann 可積分であり，
\[
\int_a^b f(x)\,dx = F(b)-F(a)
\]
が成り立つ．
\end{theorem}

\begin{proof}
$f=F'$ は連続であるから，Example 4 により Riemann 可積分である．

区間 $[a,b]$ の任意の分割
\[
\mathcal P:\ a=x_0<x_1<\cdots<x_n=b
\]
をとる．平均値の定理により，各 $i=1,\dots,n$ に対して
\[
t_i\in[x_{i-1},x_i]
\]
が存在して
\[
f(t_i)=\frac{F(x_i)-F(x_{i-1})}{x_i-x_{i-1}}
\]
となる．この点列 $t=(t_1,\dots,t_n)$ を評価点として Riemann 和をとると，
\[
\begin{aligned}
S(f,\mathcal P,t)
&=
\sum_{i=1}^n f(t_i)(x_i-x_{i-1})\\
&=
\sum_{i=1}^n \bigl(F(x_i)-F(x_{i-1})\bigr)\\
&=
F(b)-F(a)
\end{aligned}
\]
となる．

特に，$|\mathcal P_k|\to 0$ を満たす任意の分割列 $\{\mathcal P_k\}$ をとり，
各 $\mathcal P_k$ に対して上のように平均値の定理で評価点 $t^{(k)}$ を選べば，
\[
S(f,\mathcal P_k,t^{(k)})=F(b)-F(a)
\qquad (k=1,2,\dots)
\]
である．
前定理（区分求積法）より
\[
\int_a^b f(x)\,dx
=
\lim_{k\to\infty} S(f,\mathcal P_k,t^{(k)})
=
F(b)-F(a).
\]
\end{proof}


%==================================================
% 10. Remark 3 と Example 11 の差し替え
%==================================================

\begin{remark}
一様収束は，極限と積分の順序交換を保証する分かりやすい十分条件である．
しかし，一様収束よりも弱い収束では，一般には Riemann 可積分性そのものが極限で失われることがある．
次の例はその典型例である．
\end{remark}

\begin{example}
区間 $[0,1]$ 上の有理点に番号をふる：
\[
r:\mathbb N\to \mathbb Q\cap[0,1]
\]
を全単射とする．これを用いて
\[
f_n(x):=\sum_{i=1}^n \mathbf 1_{\{r(i)\}}(x)
\]
と定める．

各 $n$ に対して，$f_n$ の不連続点は有限集合
\[
\{r(1),\dots,r(n)\}
\]
に含まれる．したがって $f_n$ は区分的に連続であり，Riemann 可積分である．
実際，有限個の点の近傍だけを十分短く切り出せば
\[
\int_0^1 f_n(x)\,dx=0
\]
であることも分かる．

一方，各 $x\in[0,1]$ に対して
\[
\lim_{n\to\infty}f_n(x)=
\begin{cases}
1, & x\in \mathbb Q\cap[0,1],\\[1mm]
0, & x\in [0,1]\setminus\mathbb Q,
\end{cases}
\]
すなわち極限関数は Dirichlet 関数であり，Riemann 可積分でない．
また，$m>n$ とすると
\[
f_n(r(m))=0,\qquad \lim_{k\to\infty} f_k(r(m))=1
\]
であるから，一様収束もしていない．

したがって，一様収束よりも弱い収束では，極限関数の可積分性すら一般には保たれない．
特にこの例では
\[
\int_0^1 \lim_{n\to\infty} f_n(x)\,dx
\]
自体が定義できないので，順序交換はその前段階で破綻する．
\end{example}


%==================================================
% 11. Lebesgue 積分導入部の差し替え
%==================================================




%==================================================
% 12. 拡張実数に関する定義・注意
%==================================================

\begin{definition}[拡張実数]
\[
\overline{\mathbb R}:=\mathbb R\cup\{\pm\infty\}
\]
とおく．
\end{definition}

\begin{definition}[Darboux の上和と下和]
関数 $f:[a,b]\to\overline{\mathbb R}$ と分割
\[
\mathcal P:\ a=x_0<x_1<\cdots<x_n=b
\]
に対し，各 $i=1,\dots,n$ について
\[
M_i:=\sup_{x\in[x_{i-1},x_i]} f(x)\in\overline{\mathbb R},
\qquad
m_i:=\inf_{x\in[x_{i-1},x_i]} f(x)\in\overline{\mathbb R}
\]
とおく．
ただし，以下の有限和が $\overline{\mathbb R}$ において定義できる場合に限って，
\[
U(f,\mathcal P):=\sum_{i=1}^n M_i(x_i-x_{i-1}),
\qquad
L(f,\mathcal P):=\sum_{i=1}^n m_i(x_i-x_{i-1})
\]
をそれぞれ Darboux の上和，下和と呼ぶ．
\end{definition}

\begin{definition}[Darboux の上積分と下積分]
関数 $f:[a,b]\to\overline{\mathbb R}$ に対して
\[
\overline{\int_a^b} f(x)\,dx := \inf_{\mathcal P} U(f,\mathcal P),
\qquad
\underline{\int_a^b} f(x)\,dx := \sup_{\mathcal P} L(f,\mathcal P)
\]
と定める．
\end{definition}

\begin{definition}[Darboux 可積分]
本章では，実数値関数 $f:[a,b]\to\mathbb R$ に対して
\[
\overline{\int_a^b} f(x)\,dx
=
\underline{\int_a^b} f(x)\,dx
\in\mathbb R
\]
が成り立つとき，$f$ は Darboux 可積分であるという．
\end{definition}

\begin{remark}[拡張実数を用いるときの注意]
$\overline{\mathbb R}$ では
\[
+\infty+(-\infty)
\]
は未定義である．
したがって，ある小区間では $\sup f=+\infty$，別の小区間では $\sup f=-\infty$ となるような場合には，
$U(f,\mathcal P)$ 自体が定義できない．下和についても同様である．

また，Riemann 和
\[
S(f,\mathcal P,t)=\sum_{i=1}^n f(t_i)(x_i-x_{i-1})
\]
や Definition 3 の条件
\[
|S(f,\mathcal P,t)-I|<\varepsilon
\]
は，$f(t_i)=\pm\infty$ を許すと一般には意味を失う．
このため，本章では Darboux の上積分・下積分を述べるときに限って
$f:[a,b]\to\overline{\mathbb R}$ を許し，
Riemann 和，Riemann 可積分性，区分求積法，一様収束などを扱う箇所では
$f:[a,b]\to\mathbb R$ を仮定しておくのが安全である．

特に，Riemann 可積分な関数は自動的に有界な実数値関数である．
\end{remark}

%==================================================
% 補足
%==================================================

% Example 5 を削除するなら，
% 「Riemann 積分の線形性」
% \[
% \int_a^b (\alpha f+\beta g)\,dx
% =
% \alpha\int_a^b f(x)\,dx+\beta\int_a^b g(x)\,dx
% \]
% は独立の定理として別に置いておく方が安全です．
% Theorem 10, 11 の証明で使いやすくなります．
